\chapter{Introducción}

\section{Contexto}

El alquiler es una vía de acceso a la vivienda cada vez más importante en Lima Metropolitana. En los distritos consolidados del centro y del sur de la ciudad, la oferta de departamentos se publica sobre todo en portales inmobiliarios, donde cada aviso informa el precio pedido y un conjunto de atributos: área, número de dormitorios y de baños, antigüedad del edificio, ubicación y una descripción libre. Esa información permite estudiar con un nivel de detalle inédito cómo se forma el precio de oferta.

La literatura peruana sobre el tema ha seguido dos caminos. El primero es el enfoque hedónico clásico, que interpreta los coeficientes de una regresión del precio sobre los atributos como ``precios implícitos'' de cada característica \citep{rosen1974}. El índice hedónico de precios de departamentos del Banco Central de Reserva del Perú encuentra, por ejemplo, que el número de cocheras y el número de baños están entre las características más valoradas en Lima \citep{mundaca2018}. El segundo camino, más reciente, aplica aprendizaje automático para \emph{predecir} precios de venta o de alquiler \citep{cacnahuaray2024,barrientos2026}. Ambos son útiles, pero ninguno responde una pregunta contrafactual: \emph{¿cuánto cambiaría el precio de un departamento si tuviera un baño más, sin cambiar nada de lo demás?}

\section{Problema}

La respuesta a esa pregunta no se puede leer en una correlación ni en un coeficiente predictivo. Los departamentos con más baños son, en promedio, más grandes, más nuevos y están mejor ubicados; cada una de esas características también eleva el precio. En los datos de esta tesis, un baño adicional se asocia con un precio 45\,\% mayor cuando no se ajusta por nada, cifra que confunde el efecto del baño con el del tamaño y la zona. Al ajustar, aparece un segundo problema: ¿por qué variables ajustar? Agregar todas las disponibles no es inocuo. Algunas, como el precio por metro cuadrado, son consecuencia del precio y destruyen la estimación; otras, como los dormitorios, se deciden junto con los baños al diseñar el plano. Y queda un tercer problema que ningún ajuste resuelve: atributos no observados, como la calidad de los acabados, el piso o la vista, pueden influir tanto en el número de baños como en el precio.

El aprendizaje automático causal, en particular \textit{Double/Debiased Machine Learning} \citep{chernozhukov2018}, permite ajustar por muchas variables de forma flexible sin imponer una forma funcional y conservar una inferencia válida. Pero la flexibilidad del estimador no sustituye la estrategia de identificación \citep{hunermund2023}. La contribución de esta tesis es aplicar ese marco a un atributo residencial concreto en Lima, haciendo explícitos los supuestos, auditando los datos y midiendo cuán frágil es la conclusión.

\section{Pregunta de investigación}

\begin{quote}
¿En qué medida un diseño de aprendizaje causal con ajuste espacial permite estimar, dentro de la población de avisos comparables, la prima contrafactual sobre el precio de oferta de alquiler asociada a un baño adicional, y cuán estable es esa estimación frente a confusión no observada y a las decisiones de modelado?
\end{quote}

\section{Objetivos}

\subsection{Objetivo general}
Estimar y evaluar la robustez del efecto causal de un baño adicional sobre el precio de oferta del alquiler de departamentos en Lima Metropolitana. Para ello se usan métodos de aprendizaje causal ortogonal, controles espacialmente informados e inferencia que respeta la dependencia entre avisos.

\subsection{Objetivos específicos}
\begin{enumerate}[label=\textbf{OE\arabic*.}, leftmargin=*]
  \item Auditar, depurar y caracterizar los avisos y la selección de la muestra (procedencia, duplicados, exclusiones y perfil territorial).
  \item Definir la intervención y justificar su identificación (estimando, DAG, conjunto de ajuste, población y amenazas no identificables).
  \item Estimar y comparar la asociación hedónica con el ML causal (efecto principal con IC 95\,\%, diagnósticos de los modelos auxiliares y comparación de estimadores).
  \item Evaluar la variación socioespacial y la sensibilidad de los resultados (efectos por grupos, grilla de especificaciones, confusión no observada y dependencia espacial).
  \item Cuantificar el aporte de atributos medidos mediante IA a partir del texto de los avisos. \textit{En este entregable no se ejecuta}: la sección~\ref{sec:oe5} justifica la decisión y deja el diseño preparado.
  \item Verificar el comportamiento del procedimiento con una verdad conocida (validación semisintética: sesgo, RMSE y cobertura).
  \item Garantizar la trazabilidad de los resultados: cada tabla y figura indica el notebook y el archivo que la generan.
\end{enumerate}

\section{Hipótesis}

Las hipótesis no exigen por anticipado un efecto significativo ni positivo:
\begin{description}[leftmargin=1.2cm]
  \item[H1.] La magnitud de la asociación inicial cambia de forma mensurable tras el ajuste preespecificado.
  \item[H2.] La omisión de confusores o la inclusión de controles causalmente inadecuados altera la estimación más que la elección del algoritmo de predicción.
  \item[H3.] Existe heterogeneidad del efecto según el contexto territorial o el tamaño del departamento, y solo se afirmará si la precisión estadística la respalda.
  \item[H4.] La solidez de la conclusión depende de la fuerza que se postule para los confusores no observados.
\end{description}

\section{Alcance y aportes}

\textbf{Alcance.} El objeto son los precios \emph{de oferta} mensuales de departamentos publicados en un portal, no contratos firmados ni el bienestar de los hogares. La cobertura son 13 distritos de Lima Metropolitana con avisos activos en marzo de 2025. La unidad de análisis es el aviso, con un manejo explícito de los avisos repetidos y de la agrupación por edificio y por zona. Los resultados no se extrapolan al parque residencial completo ni a otros distritos.

\textbf{Aportes verificables.} (i) Una auditoría reproducible de la base, que documenta duplicados entre distritos y errores de geocodificación que no se habían detectado. (ii) Una estimación del efecto de un baño adicional con identificación explícita, inferencia agrupada y comparación con el enfoque hedónico. (iii) Una demostración empírica de cómo la elección de controles cambia el resultado. (iv) Un análisis de sensibilidad calibrado con confusores observados. (v) Una validación semisintética anclada en las covariables reales del mercado limeño.

\textbf{Limitaciones principales.} El tratamiento no se asignó al azar, por lo que la interpretación causal depende de la ausencia de confusores omitidos importantes, supuesto que solo puede acotarse. La muestra corresponde a un único corte de un portal, y la antigüedad falta en el 25\,\% de la muestra final. Los atributos de calidad del texto aún no se extraen (OE5).

\section{Organización del documento}

El capítulo~\ref{cap:estado} revisa la literatura. El capítulo~\ref{cap:marco} presenta el modelo causal, el estimando y el DAG. El capítulo~\ref{cap:datos} describe los datos, la auditoría y los métodos de estimación, y el capítulo~\ref{cap:diseno} el diseño de las pruebas de robustez y de la validación semisintética. El capítulo~\ref{cap:resultados} reporta los resultados, el~\ref{cap:discusion} los discute y el~\ref{cap:conclusiones} presenta las conclusiones por objetivo. Los anexos contienen el diccionario de variables, el registro de la búsqueda bibliográfica, el inventario de artefactos y las versiones de software.
